\renewcommand{\tema}{Preguntas Post-Clase - Fuerzas y Leyes de Newton}

\twocolumn
\section{PREGUNTAS POST-CLASE}

\begin{preguntabox}
\subsection{Pregunta 1}
¿Cuál de las siguientes situaciones ilustra mejor la Primera Ley de Newton?

\begin{enumerate}
    \item[A)] Un automóvil acelera desde el reposo
    \item[B)] Un libro permanece en reposo sobre una mesa sin fuerzas externas
    \item[C)] Una pelota cae verticalmente bajo la acción de la gravedad
    \item[D)] Un disco de hockey sobre hielo continúa deslizándose con velocidad casi constante
\end{enumerate}
\end{preguntabox}

\begin{preguntabox}
\subsection{Pregunta 2}
¿Cuál es la aceleración que experimenta un cuerpo de 5 kg cuando se le aplica una fuerza neta de 20 N?

\begin{enumerate}
    \item[A)] 4 m/s$^2$
    \item[B)] 25 m/s$^2$
    \item[C)] 100 m/s$^2$
    \item[D)] 0.25 m/s$^2$
\end{enumerate}
\end{preguntabox}

\begin{preguntabox}
\subsection{Pregunta 3}
Si la masa de un objeto se duplica mientras la fuerza neta aplicada permanece constante, ¿qué sucede con su aceleración?

\begin{enumerate}
    \item[A)] Se duplica
    \item[B)] Se cuadruplica
    \item[C)] Se reduce a la mitad
    \item[D)] Permanece constante
\end{enumerate}
\end{preguntabox}

\begin{preguntabox}
\subsection{Pregunta 4}
Un libro descansa sobre una mesa. La mesa ejerce una fuerza hacia arriba sobre el libro. Según la Tercera Ley de Newton, ¿cuál es la reacción a esta fuerza?

\begin{enumerate}
    \item[A)] El peso del libro
    \item[B)] La fuerza que el libro ejerce hacia abajo sobre la mesa
    \item[C)] La fuerza que la Tierra ejerce sobre el libro
    \item[D)] La fuerza de fricción entre el libro y la mesa
\end{enumerate}
\end{preguntabox}

\begin{preguntabox}
\subsection{Pregunta 5}
¿Cuál es el peso de un objeto de 8 kg en la superficie terrestre? (Usa $g = 10$ m/s$^2$)

\begin{enumerate}
    \item[A)] 8 N
    \item[B)] 18 N
    \item[C)] 80 N
    \item[D)] 800 N
\end{enumerate}
\end{preguntabox}

\begin{preguntabox}
\subsection{Pregunta 6}
Un bloque de 10 kg es empujado sobre una superficie horizontal sin fricción con una fuerza de 40 N. Si el bloque parte del reposo, ¿qué velocidad alcanza después de 5 segundos?

\begin{enumerate}
    \item[A)] 4 m/s
    \item[B)] 8 m/s
    \item[C)] 20 m/s
    \item[D)] 200 m/s
\end{enumerate}
\end{preguntabox}

\begin{preguntabox}
\subsection{Pregunta 7}
Una caja de 20 kg cuelga del techo mediante una cuerda. Si la caja está en reposo, ¿cuál es la tensión en la cuerda? (Usa $g = 10$ m/s$^2$)

\begin{enumerate}
    \item[A)] 20 N
    \item[B)] 100 N
    \item[C)] 200 N
    \item[D)] 0 N
\end{enumerate}
\end{preguntabox}

\begin{preguntabox}
\subsection{Pregunta 8}
Un nadador empuja la pared de una alberca con una fuerza de 150 N. Según la Tercera Ley de Newton, ¿qué fuerza ejerce la pared sobre el nadador?

\begin{enumerate}
    \item[A)] 0 N
    \item[B)] 75 N en la misma dirección
    \item[C)] 150 N en dirección opuesta
    \item[D)] 300 N en dirección opuesta
\end{enumerate}
\end{preguntabox}

\begin{preguntabox}
\subsection{Pregunta 9}
Una viga horizontal de 4 m está apoyada en su punto medio. Se colocan pesos de 30 N a 1 m del pivote y de masa desconocida a 2 m del pivote en el lado opuesto. Para que la viga esté en equilibrio rotacional, ¿cuál debe ser el peso desconocido?

\begin{enumerate}
    \item[A)] 15 N
    \item[B)] 30 N
    \item[C)] 60 N
    \item[D)] 120 N
\end{enumerate}
\end{preguntabox}

\begin{preguntabox}
\subsection{Pregunta 10}
Dos resortes idénticos, cada uno con constante $k = 200$ N/m, se conectan en serie (uno después del otro). Si se aplica una fuerza de 40 N al sistema, ¿cuál es la elongación total?

\begin{enumerate}
    \item[A)] 0.1 m
    \item[B)] 0.2 m
    \item[C)] 0.4 m
    \item[D)] 0.8 m
\end{enumerate}
\end{preguntabox}

\begin{preguntabox}
\subsection{Pregunta 11}
Un bloque de 5 kg se encuentra sobre una mesa. Se le aplica una fuerza horizontal de 30 N, pero el bloque no se mueve debido a la fricción estática. ¿Cuál de las siguientes afirmaciones es correcta?

\begin{enumerate}
    \item[A)] La fuerza neta sobre el bloque es 30 N
    \item[B)] La fuerza de fricción estática es 30 N en dirección opuesta a la fuerza aplicada
    \item[C)] El bloque está acelerando lentamente
    \item[D)] No hay fuerzas actuando sobre el bloque
\end{enumerate}
\end{preguntabox}

\begin{preguntabox}
\subsection{Pregunta 12}
Una fuerza de 50 N actúa sobre un bloque formando un ángulo de 60$^\circ$ con la horizontal. ¿Cuál es la componente horizontal de esta fuerza? (Usa cos 60$^\circ$ = 0.5)

\begin{enumerate}
    \item[A)] 25 N
    \item[B)] 43.3 N
    \item[C)] 50 N
    \item[D)] 86.6 N
\end{enumerate}
\end{preguntabox}

\begin{preguntabox}
\subsection{Pregunta 13}
Dos masas de 4 kg y 6 kg están separadas por una distancia de 3 m. Si la distancia se aumenta a 6 m, ¿cómo cambia la fuerza gravitacional entre ellas?

\begin{enumerate}
    \item[A)] Se reduce a la mitad
    \item[B)] Se reduce a un cuarto
    \item[C)] Se duplica
    \item[D)] Se cuadruplica
\end{enumerate}
\end{preguntabox}

\begin{preguntabox}
\subsection{Pregunta 14}
Un planeta X tiene el doble del radio de la Tierra pero la misma masa. ¿Cómo se compara la aceleración gravitacional en la superficie del planeta X ($g_X$) con la de la Tierra ($g_T$)?

\begin{enumerate}
    \item[A)] $g_X = g_T$
    \item[B)] $g_X = 2g_T$
    \item[C)] $g_X = g_T/2$
    \item[D)] $g_X = g_T/4$
\end{enumerate}
\end{preguntabox}

\begin{preguntabox}
\subsection{Pregunta 15}
Un planeta tiene un período orbital de 8 años terrestres alrededor del Sol. Si otro planeta tiene un período orbital de 27 años, ¿cuál es la relación entre sus distancias al Sol? (distancia del primer planeta / distancia del segundo planeta)

\begin{enumerate}
    \item[A)] 2/3
    \item[B)] 4/9
    \item[C)] 8/27
    \item[D)] 64/729
\end{enumerate}
\end{preguntabox}

\newpage
\onecolumn
\section{RESPUESTAS Y RETROALIMENTACIÓN}

\begin{respuestabox}
\subsection{Pregunta 1}
\textcolor{verdecorrecto}{\textbf{Respuesta correcta: D) Un disco de hockey sobre hielo continúa deslizándose con velocidad casi constante}}

\textbf{Justificación:}\\
\textbf{Concepto:} Primera Ley de Newton (Lección: Sección 3.1)

La Primera Ley establece que un objeto mantiene su estado de movimiento (reposo o velocidad constante) cuando la fuerza neta es cero. El disco de hockey experimenta muy poca fricción sobre el hielo, por lo que la fuerza neta es aproximadamente cero y continúa moviéndose con velocidad casi constante.

\textbf{Condición:}
$$\sum \vec{F} \approx 0 \quad \Rightarrow \quad \vec{v} = \text{constante}$$

\textbf{Respuestas incorrectas:}
\begin{itemize}
    \item \textbf{A):} Hay fuerza neta (del motor), por lo que hay aceleración. Ilustra la Segunda Ley, no la Primera.
    \item \textbf{B):} Incorrecto decir "sin fuerzas externas". Hay peso y normal que se equilibran.
    \item \textbf{C):} Hay fuerza neta (peso), por lo que acelera. No ilustra la Primera Ley.
\end{itemize}
\end{respuestabox}

\begin{respuestabox}
\subsection{Pregunta 2}
\textcolor{verdecorrecto}{\textbf{Respuesta correcta: A) 4 m/s$^2$}}

\textbf{Justificación:}\\
\textbf{Fórmula utilizada:} Segunda Ley de Newton (Lección: Sección 3.2)
$$F_{neta} = ma$$

\textbf{Despeje para $a$:}
$$a = \frac{F_{neta}}{m}$$

\textbf{Sustitución:}
$$a = \frac{20 \text{ N}}{5 \text{ kg}} = 4 \text{ m/s}^2$$

\textbf{Respuestas incorrectas:}
\begin{itemize}
    \item \textbf{B):} Error aritmético. Sumó en lugar de dividir: $20 + 5 = 25$.
    \item \textbf{C):} Error aritmético. Multiplicó en lugar de dividir: $20 \times 5 = 100$.
    \item \textbf{D):} Error aritmético. Invirtió la división: $5/20 = 0.25$.
\end{itemize}
\end{respuestabox}

\begin{respuestabox}
\subsection{Pregunta 3}
\textcolor{verdecorrecto}{\textbf{Respuesta correcta: C) Se reduce a la mitad}}

\textbf{Justificación:}\\
\textbf{Fórmula utilizada:} Segunda Ley de Newton (Lección: Sección 3.2)
$$a = \frac{F}{m}$$

La aceleración es inversamente proporcional a la masa cuando la fuerza es constante.

Si la masa se duplica ($m' = 2m$):
$$a' = \frac{F}{2m} = \frac{1}{2} \cdot \frac{F}{m} = \frac{a}{2}$$

Por lo tanto, la aceleración se reduce a la mitad.

\textbf{Respuestas incorrectas:}
\begin{itemize}
    \item \textbf{A):} Error conceptual. Confunde relación directa con inversa.
    \item \textbf{B):} Error conceptual. Aplicaría si hubiera una relación cuadrática.
    \item \textbf{D):} Error conceptual. Ignora la dependencia de $a$ con $m$.
\end{itemize}
\end{respuestabox}

\begin{respuestabox}
\subsection{Pregunta 4}
\textcolor{verdecorrecto}{\textbf{Respuesta correcta: B) La fuerza que el libro ejerce hacia abajo sobre la mesa}}

\textbf{Justificación:}\\
La Tercera Ley de Newton (Lección: Sección 3.5) establece que para cada acción hay una reacción igual en magnitud, opuesta en dirección y que actúa sobre cuerpos diferentes.

Si la mesa ejerce una fuerza $\vec{F}$ hacia arriba sobre el libro (fuerza normal), el par acción-reacción es la fuerza que el libro ejerce hacia abajo sobre la mesa con magnitud $\vec{F}$.

\textbf{Criterios para identificar pares acción-reacción:}
- Misma magnitud
- Direcciones opuestas
- Actúan sobre cuerpos diferentes
- Misma naturaleza (ambas son fuerzas de contacto)

\textbf{Respuestas incorrectas:}
\begin{itemize}
    \item \textbf{A):} El peso actúa sobre el libro, no sobre la mesa. Peso y normal no forman par acción-reacción.
    \item \textbf{C):} Es el peso, no el par de la normal. Actúa sobre el libro, no sobre la mesa.
    \item \textbf{D):} Irrelevante para este caso. No es el par de la fuerza normal.
\end{itemize}
\end{respuestabox}

\begin{respuestabox}
\subsection{Pregunta 5}
\textcolor{verdecorrecto}{\textbf{Respuesta correcta: C) 80 N}}

\textbf{Justificación:}\\
\textbf{Fórmula utilizada:} Peso (Lección: Sección 3.4)
$$W = mg$$

donde:
- $W$ = peso (N)
- $m$ = masa (kg)
- $g$ = aceleración gravitacional (m/s$^2$)

\textbf{Sustitución:}
$$W = (8 \text{ kg})(10 \text{ m/s}^2) = 80 \text{ N}$$

\textbf{Respuestas incorrectas:}
\begin{itemize}
    \item \textbf{A):} Error conceptual. Confunde masa con peso, usando solo el valor numérico de la masa.
    \item \textbf{B):} Error aritmético, suma en lugar de multiplicación: $8 + 10 = 18$.
    \item \textbf{D):} Error de unidades o de cálculo. Multiplicación incorrecta con un cero extra.
\end{itemize}
\end{respuestabox}

\begin{respuestabox}
\subsection{Pregunta 6}
\textcolor{verdecorrecto}{\textbf{Respuesta correcta: C) 20 m/s}}

\textbf{Justificación:}\\
\textbf{Fórmulas utilizadas:} Segunda Ley de Newton y cinemática (Lección: Sección 3.2)

\textbf{Paso 1:} Calcular la aceleración
$$a = \frac{F}{m} = \frac{40 \text{ N}}{10 \text{ kg}} = 4 \text{ m/s}^2$$

\textbf{Paso 2:} Calcular la velocidad final (parte del reposo: $v_0 = 0$)
$$v = v_0 + at = 0 + (4 \text{ m/s}^2)(5 \text{ s}) = 20 \text{ m/s}$$

\textbf{Respuestas incorrectas:}
\begin{itemize}
    \item \textbf{A):} Error conceptual. Confunde aceleración con velocidad final.
    \item \textbf{B):} Error aritmético en el cálculo.
    \item \textbf{D):} Error conceptual. Multiplicó fuerza por tiempo (confunde impulso con velocidad).
\end{itemize}
\end{respuestabox}

\begin{respuestabox}
\subsection{Pregunta 7}
\textcolor{verdecorrecto}{\textbf{Respuesta correcta: C) 200 N}}

\textbf{Justificación:}\\
\textbf{Concepto:} Equilibrio traslacional (Lección: Sección 4.1)

La caja está en reposo, por lo que está en equilibrio traslacional:
$$\sum F_y = 0$$

\textbf{Fuerzas en dirección vertical:}
- Tensión $T$ (hacia arriba, +)
- Peso $W$ (hacia abajo, -)

\textbf{Aplicando la condición de equilibrio:}
$$T - W = 0 \quad \Rightarrow \quad T = W$$

\textbf{Cálculo del peso:}
$$W = mg = (20 \text{ kg})(10 \text{ m/s}^2) = 200 \text{ N}$$

Por lo tanto: $T = 200$ N

\textbf{Respuestas incorrectas:}
\begin{itemize}
    \item \textbf{A):} Error conceptual. Confunde masa con peso.
    \item \textbf{B):} Error aritmético. Usó $g = 5$ m/s$^2$ incorrectamente.
    \item \textbf{D):} Error conceptual grave. Ignora que debe haber tensión para equilibrar el peso.
\end{itemize}
\end{respuestabox}

\begin{respuestabox}
\subsection{Pregunta 8}
\textcolor{verdecorrecto}{\textbf{Respuesta correcta: C) 150 N en dirección opuesta}}

\textbf{Justificación:}\\
\textbf{Concepto:} Tercera Ley de Newton (Lección: Sección 3.5)
$$\vec{F}_{AB} = -\vec{F}_{BA}$$

Si el nadador ejerce 150 N sobre la pared, por la Tercera Ley, la pared ejerce 150 N sobre el nadador en dirección opuesta (empujando al nadador hacia atrás).

\textbf{Características del par acción-reacción:}
\begin{itemize}
    \item Misma magnitud: 150 N
    \item Direcciones opuestas
    \item Actúan sobre cuerpos diferentes (nadador y pared)
    \item Son simultáneas
\end{itemize}

\textbf{Respuestas incorrectas:}
\begin{itemize}
    \item \textbf{A):} Error conceptual. Viola la Tercera Ley de Newton.
    \item \textbf{B):} Error conceptual doble. Magnitud incorrecta y misma dirección.
    \item \textbf{D):} Error conceptual. La magnitud debe ser igual, no el doble.
\end{itemize}
\end{respuestabox}

\begin{respuestabox}
\subsection{Pregunta 9}
\textcolor{verdecorrecto}{\textbf{Respuesta correcta: A) 15 N}}

\textbf{Justificación:}\\
\textbf{Concepto:} Segunda condición de equilibrio - equilibrio rotacional (Lección: Sección 4.3)
$$\sum \tau = 0$$

\textbf{Configuración:}
- Peso 1: $F_1 = 30$ N a distancia $r_1 = 1$ m
- Peso 2: $F_2 = ?$ a distancia $r_2 = 2$ m

\textbf{Torcas respecto al pivote:}
$$\tau_1 = r_1 F_1 = (1 \text{ m})(30 \text{ N}) = 30 \text{ N·m (horaria, -)}$$
$$\tau_2 = r_2 F_2 = (2 \text{ m})(F_2) \text{ (antihoraria, +)}$$

\textbf{Condición de equilibrio:}
$$\tau_2 - \tau_1 = 0$$
$$(2)(F_2) = 30$$
$$F_2 = 15 \text{ N}$$

\textbf{Respuestas incorrectas:}
\begin{itemize}
    \item \textbf{B):} Error conceptual. Asume que los pesos deben ser iguales, ignorando las distancias diferentes.
    \item \textbf{C):} Error aritmético. Multiplicó en lugar de dividir: $30 \times 2 = 60$.
    \item \textbf{D):} Error de cálculo múltiple.
\end{itemize}
\end{respuestabox}

\begin{respuestabox}
\subsection{Pregunta 10}
\textcolor{verdecorrecto}{\textbf{Respuesta correcta: C) 0.4 m}}

\textbf{Justificación:}\\
\textbf{Concepto:} Ley de Hooke para resortes en serie (Lección: Sección 5)

Cuando dos resortes idénticos están en serie, la constante efectiva es:
$$k_{efectiva} = \frac{k}{2} = \frac{200 \text{ N/m}}{2} = 100 \text{ N/m}$$

\textbf{Fórmula:}
$$F = k_{efectiva} \cdot x_{total}$$

\textbf{Despeje:}
$$x_{total} = \frac{F}{k_{efectiva}} = \frac{40 \text{ N}}{100 \text{ N/m}} = 0.4 \text{ m}$$

\textbf{Explicación física:} En serie, cada resorte experimenta la misma fuerza (40 N), por lo que cada uno se estira:
$$x_{individual} = \frac{40}{200} = 0.2 \text{ m}$$

La elongación total es: $x_{total} = 0.2 + 0.2 = 0.4$ m

\textbf{Respuestas incorrectas:}
\begin{itemize}
    \item \textbf{A):} Usó $k_{efectiva} = 2k = 400$ N/m (fórmula de paralelo, no serie).
    \item \textbf{B):} Calculó la elongación de un solo resorte, olvidando que hay dos.
    \item \textbf{D):} Error de cálculo, duplicó incorrectamente.
\end{itemize}
\end{respuestabox}

\begin{respuestabox}
\subsection{Pregunta 11}
\textcolor{verdecorrecto}{\textbf{Respuesta correcta: B) La fuerza de fricción estática es 30 N en dirección opuesta a la fuerza aplicada}}

\textbf{Justificación:}\\
\textbf{Concepto:} Primera condición de equilibrio y fricción estática (Lección: Sección 4.1)

Si el bloque no se mueve, está en equilibrio traslacional:
$$\sum F_x = 0$$

\textbf{Fuerzas horizontales:}
- Fuerza aplicada: $F_{aplicada} = 30$ N (hacia la derecha, +)
- Fuerza de fricción estática: $f_s$ (hacia la izquierda, -)

\textbf{Condición de equilibrio:}
$$F_{aplicada} - f_s = 0$$
$$f_s = 30 \text{ N}$$

La fricción se ajusta automáticamente para igualar la fuerza aplicada (hasta su valor máximo).

\textbf{Respuestas incorrectas:}
\begin{itemize}
    \item \textbf{A):} Error. Si la fuerza neta fuera 30 N, el bloque aceleraría.
    \item \textbf{C):} Error. Si no se mueve, su aceleración es cero.
    \item \textbf{D):} Error grave. Hay múltiples fuerzas, pero están equilibradas.
\end{itemize}
\end{respuestabox}

\begin{respuestabox}
\subsection{Pregunta 12}
\textcolor{verdecorrecto}{\textbf{Respuesta correcta: A) 25 N}}

\textbf{Justificación:}\\
\textbf{Concepto:} Componentes de un vector (Lección: Sección 1.2)

\textbf{Fórmula:}
$$F_x = F \cos\theta$$

\textbf{Datos:}
- $F = 50$ N
- $\theta = 60^\circ$
- $\cos 60^\circ = 0.5$

\textbf{Sustitución:}
$$F_x = 50 \text{ N} \times 0.5 = 25 \text{ N}$$

\textbf{Respuestas incorrectas:}
\begin{itemize}
    \item \textbf{B):} Error conceptual. Usó seno en lugar de coseno: $F_y = 50 \sin 60^\circ = 50 \times 0.866 \approx 43.3$ N (componente vertical).
    \item \textbf{C):} Error conceptual. Confunde la magnitud total con la componente.
    \item \textbf{D):} Error de cálculo. Duplicó incorrectamente la componente vertical.
\end{itemize}
\end{respuestabox}

\begin{respuestabox}
\subsection{Pregunta 13}
\textcolor{verdecorrecto}{\textbf{Respuesta correcta: B) Se reduce a un cuarto}}

\textbf{Justificación:}\\
\textbf{Fórmula utilizada:} Ley de Gravitación Universal (Lección: Sección 6)
$$F_g = G\frac{m_1 m_2}{r^2}$$

La fuerza es inversamente proporcional al cuadrado de la distancia.

\textbf{Situación inicial:} $r_i = 3$ m
$$F_i = G\frac{m_1 m_2}{(3)^2} = G\frac{m_1 m_2}{9}$$

\textbf{Situación final:} $r_f = 6$ m
$$F_f = G\frac{m_1 m_2}{(6)^2} = G\frac{m_1 m_2}{36}$$

\textbf{Relación:}
$$\frac{F_f}{F_i} = \frac{9}{36} = \frac{1}{4}$$

Por lo tanto, la fuerza se reduce a un cuarto (o se divide entre 4).

\textbf{Respuestas incorrectas:}
\begin{itemize}
    \item \textbf{A):} Error conceptual. Ignora el exponente cuadrado en la dependencia con $r$.
    \item \textbf{C):} Error conceptual grave. Confunde aumento de distancia con aumento de fuerza.
    \item \textbf{D):} Error conceptual. Invierte la relación.
\end{itemize}
\end{respuestabox}

\begin{respuestabox}
\subsection{Pregunta 14}
\textcolor{verdecorrecto}{\textbf{Respuesta correcta: D) $g_X = g_T/4$}}

\textbf{Justificación:}\\
\textbf{Fórmula utilizada:} Aceleración gravitacional en la superficie (Lección: Sección 6.1)
$$g = G\frac{M}{R^2}$$

\textbf{Para la Tierra:}
$$g_T = G\frac{M_T}{R_T^2}$$

\textbf{Para el planeta X:} (masa igual, radio doble)
$$g_X = G\frac{M_X}{R_X^2} = G\frac{M_T}{(2R_T)^2} = G\frac{M_T}{4R_T^2} = \frac{1}{4} \cdot G\frac{M_T}{R_T^2} = \frac{g_T}{4}$$

Por lo tanto: $g_X = g_T/4$

\textbf{Interpretación:} Un planeta con el doble de radio pero la misma masa tiene una aceleración gravitacional cuatro veces menor en su superficie.

\textbf{Respuestas incorrectas:}
\begin{itemize}
    \item \textbf{A):} Error. Ignora el efecto del radio en la aceleración gravitacional.
    \item \textbf{B):} Error conceptual. Confunde relación directa con inversa cuadrática.
    \item \textbf{C):} Error. No considera el exponente cuadrado del radio.
\end{itemize}
\end{respuestabox}

\begin{respuestabox}
\subsection{Pregunta 15}
\textcolor{verdecorrecto}{\textbf{Respuesta correcta: B) 4/9}}

\textbf{Justificación:}\\
\textbf{Fórmula utilizada:} Tercera Ley de Kepler (Lección: Sección 7.3)
$$\frac{T_1^2}{T_2^2} = \frac{a_1^3}{a_2^3}$$

\textbf{Datos:}
- $T_1 = 8$ años
- $T_2 = 27$ años

\textbf{Aplicando la Tercera Ley:}
$$\frac{a_1^3}{a_2^3} = \frac{T_1^2}{T_2^2} = \frac{8^2}{27^2} = \frac{64}{729}$$

\textbf{Para encontrar la relación de distancias:}
$$\frac{a_1}{a_2} = \sqrt[3]{\frac{64}{729}} = \frac{\sqrt[3]{64}}{\sqrt[3]{729}} = \frac{4}{9}$$

\textbf{Verificación:}
- $\sqrt[3]{64} = 4$ (porque $4^3 = 64$)
- $\sqrt[3]{729} = 9$ (porque $9^3 = 729$)

\textbf{Respuestas incorrectas:}
\begin{itemize}
    \item \textbf{A):} Error de cálculo. Simplificación incorrecta.
    \item \textbf{C):} Error conceptual. Usó la relación de períodos directamente sin elevar al cuadrado ni sacar raíz cúbica.
    \item \textbf{D):} Error conceptual. Calculó $T_1^2/T_2^2$ pero olvidó sacar la raíz cúbica para obtener $a_1/a_2$.
\end{itemize}
\end{respuestabox}